\subsection{Ambiente de Trabajo [Pendientes: 1. acordar plataforma grupal y mecanismo de seguimiento de acuerdos]}\label{sec:spmp-12-1}
El equipo adopta como base de convivencia el canvas de trabajo de grado aportado, cuyo propósito es aplicar los conocimientos adquiridos en la carrera y unir esfuerzos para desarrollar un trabajo de calidad. Sus valores son responsabilidad, compromiso, sinceridad, comunicación y empatía, y sus prioridades sitúan la calidad sobre la rapidez, la comunicación sobre la imposición y la responsabilidad sobre la procrastinación. La iniciativa, el reconocimiento de errores y los aportes de cada integrante deben acompañarse de comunicación asertiva y consideración por el trabajo común, de manera que una dificultad individual pueda tratarse oportunamente como un problema que el equipo ayuda a resolver~\cite{canvas2026}.

Las decisiones se tomarán mediante diálogo y consenso, y cuando persista el desacuerdo se recurrirá a votación. En caso de empate, el grupo debatirá para argumentar las posturas frente a validación externa, según la regla expresada en el canvas, sin atribuir de antemano esa validación a una persona que no haya aceptado intervenir. Las reuniones serán presenciales o virtuales, previamente acordadas y con un objetivo que permita revisar avances, dificultades o entregas. También se establecerán plazos internos anteriores a los institucionales para revisar y corregir, respetando los horarios, compromisos y actividades individuales. La organización formal del semestre deberá concretar la agenda y las herramientas para distribuir el trabajo.

La retroalimentación será respetuosa, reconocerá aspectos positivos e identificará oportunidades de mejora, buscando soluciones conjuntas y asumiendo el error como una posibilidad de aprendizaje colectivo. El canvas también propone espacios de convivencia y confianza para conocer habilidades y facilitar la participación de todos. Para aplicar estas reglas, se propone conservar acuerdos y dificultades relevantes junto con su responsable de seguimiento, utilizar la plataforma grupal que se acuerde y revisar en las reuniones si las acciones comprometidas se realizaron. Estos mecanismos requieren ratificación y asignación operativa, mientras que los valores y reglas anteriores corresponden al material existente del equipo.
